\section{Connecting to a Server}

\subsection{What is SSH?}

SSH (Secure Shell) is the standard method for connecting to a remote Linux machine.

\begin{lstlisting}
ssh username@pelee.math.private.uwaterloo.ca
\end{lstlisting}

\begin{lstlisting}
ssh username@rondeau.math.private.uwaterloo.ca
\end{lstlisting}

\subsection{Windows}
Open Windows Terminal and run the SSH command.

\subsection{macOS and Linux}
Open a terminal and run the SSH command.

\textbf{Note:} Nothing appears on the screen while typing your password. This is normal.

\subsection{First Login}

Type \texttt{yes} when asked whether you trust the host.

\subsection{Successful Login}

\begin{lstlisting}
username@pelee:~$
\end{lstlisting}

\subsection{Connecting from Off Campus (VPN)}

Pelee and Rondeau are on the campus network, so from off campus you must connect to the
university VPN first. If \texttt{ssh} hangs and then times out, the most common reason is that the
VPN is not connected.

The university uses the Cisco Secure Client. Installation differs by operating system and changes
from time to time, so follow the current instructions from IST:

\begin{center}
\url{https://uwaterloo.atlassian.net/wiki/spaces/CWKB/pages/1736840329}
\end{center}

\begin{itemize}
\item Connect to the VPN \emph{before} running \texttt{ssh}.
\item In the second-password field, enter your Duo method; \texttt{push} sends a notification to
the Duo Mobile app.
\item You do \emph{not} need the VPN for the Alliance clusters (Appendix~\ref{app:alliance}).
\end{itemize}

\subsection{Logging Out}

\begin{lstlisting}
exit
\end{lstlisting}

\subsection{Exercise}

\begin{enumerate}
\item Connect to Pelee.
\item Run \texttt{pwd}.
\item Disconnect using \texttt{exit}.
\end{enumerate}
